% ═══════════════════════════════════════════════════════════════
% MUSTERLÖSUNG: Tabellenkalkulation Grundlagen
% Informatik Klasse 7 | Stunde 3
% ═══════════════════════════════════════════════════════════════

\documentclass[11pt, a4paper]{article}

\usepackage[utf8]{inputenc}
\usepackage[T1]{fontenc}
\usepackage[ngerman]{babel}

\usepackage{helvet}
\renewcommand{\familydefault}{\sfdefault}

\usepackage[margin=1.5cm, top=1.8cm, bottom=1.5cm]{geometry}
\usepackage{enumitem}
\usepackage{tabularx}
\usepackage{booktabs}
\usepackage{array}
\usepackage{tikz}
\usepackage{fancyhdr}
\usepackage{xcolor}
\usepackage{tcolorbox}
\tcbuselibrary{skins}
\usepackage{amssymb}

% ═══════════════════════════════════════════════════════════════
% FARBEN
% ═══════════════════════════════════════════════════════════════

\definecolor{informatik}{HTML}{b35c00}
\definecolor{hellgrau}{HTML}{F5F5F5}
\definecolor{mittelgrau}{HTML}{888888}
\definecolor{loesung}{HTML}{c62828}

\colorlet{fachfarbe}{informatik}

% ═══════════════════════════════════════════════════════════════
% KOPFZEILE
% ═══════════════════════════════════════════════════════════════

\pagestyle{fancy}
\fancyhf{}
\renewcommand{\headrulewidth}{0.4pt}
\renewcommand{\footrulewidth}{0pt}
\lhead{\small Informatik Klasse 7 | Tabellenkalkulation}
\rhead{\small \textcolor{loesung}{\textbf{MUSTERLÖSUNG}}}

% ═══════════════════════════════════════════════════════════════
% BOXEN
% ═══════════════════════════════════════════════════════════════

\newtcolorbox{merkkasten}[1][]{
    colback=hellgrau,
    colframe=hellgrau,
    boxrule=0pt,
    arc=0pt,
    left=8pt, right=8pt, top=8pt, bottom=6pt,
    borderline north={2pt}{0pt}{fachfarbe},
    before skip=8pt, after skip=8pt,
    #1
}

\newtcolorbox{lueckenkasten}[1][]{
    colback=white,
    colframe=fachfarbe,
    boxrule=1pt,
    arc=0pt,
    left=8pt, right=8pt, top=8pt, bottom=8pt,
    before skip=6pt, after skip=6pt,
    #1
}

% Lösungstext
\newcommand{\lsg}[1]{\textcolor{loesung}{\textbf{#1}}}

\setlength{\parindent}{0pt}
\setlength{\parskip}{4pt}

% ═══════════════════════════════════════════════════════════════
% DOKUMENT
% ═══════════════════════════════════════════════════════════════

\begin{document}

% Header
\begin{center}
    {\Large\bfseries\textcolor{fachfarbe}{Tabellenkalkulation Grundlagen}}\\[0.3em]
    {\small Arbeitsblatt | \textcolor{loesung}{MUSTERLÖSUNG}}
\end{center}
\vspace{-0.3em}
\hrule
\vspace{0.6em}

% ═══════════════════════════════════════════════════════════════
% KASTEN 1: Zelle
% ═══════════════════════════════════════════════════════════════

\textbf{\textcolor{fachfarbe}{Kasten 1: Was ist eine Zelle?}}

\begin{lueckenkasten}
Eine \textbf{Zelle} ist \lsg{ein einzelnes Kästchen,}

\vspace{0.3em}
in das man \lsg{Daten} eingeben kann.
\end{lueckenkasten}

% ═══════════════════════════════════════════════════════════════
% KASTEN 2: Spalten und Zeilen
% ═══════════════════════════════════════════════════════════════

\textbf{\textcolor{fachfarbe}{Kasten 2: Spalten und Zeilen}}

\begin{lueckenkasten}
\begin{tabular}{@{}l@{\hspace{1cm}}l@{}}
\textbf{Spalte} & verläuft \lsg{senkrecht} (↓) \\[0.5em]
                & wird mit \lsg{Buchstaben} bezeichnet (A, B, C, ...) \\[0.8em]
\textbf{Zeile}  & verläuft \lsg{waagerecht} (→) \\[0.5em]
                & wird mit \lsg{Nummern} bezeichnet (1, 2, 3, ...)
\end{tabular}

\vspace{0.5em}
\textit{Merkhilfe:} \textbf{S}palte = \textbf{S}\lsg{enkrecht}
\end{lueckenkasten}

% ═══════════════════════════════════════════════════════════════
% KASTEN 3: Formeln
% ═══════════════════════════════════════════════════════════════

\textbf{\textcolor{fachfarbe}{Kasten 3: Formeln eingeben}}

\begin{lueckenkasten}
Formeln beginnen \textbf{immer} mit dem Zeichen: \lsg{=}

\vspace{0.3em}
Beispiel: \texttt{\lsg{=}B2*C2}
\end{lueckenkasten}

% ═══════════════════════════════════════════════════════════════
% KASTEN 4: Rechenzeichen
% ═══════════════════════════════════════════════════════════════

\textbf{\textcolor{fachfarbe}{Kasten 4: Rechenzeichen in Formeln}}

\begin{lueckenkasten}
\begin{center}
\begin{tabular}{c c c}
\toprule
\textbf{Mathe} & \textbf{Taste} & \textbf{Bedeutung} \\
\midrule
$+$ & \texttt{+} & Addition \\[0.3em]
$-$ & \texttt{-} & Subtraktion \\[0.3em]
$\times$ & \lsg{*} & Multiplikation \\[0.3em]
$\div$ & \lsg{/} & Division \\
\bottomrule
\end{tabular}
\end{center}
\end{lueckenkasten}

% ═══════════════════════════════════════════════════════════════
% KASTEN 5: SUMME-Funktion
% ═══════════════════════════════════════════════════════════════

\textbf{\textcolor{fachfarbe}{Kasten 5: Die SUMME-Funktion}}

\begin{lueckenkasten}
Die Formel \texttt{=SUMME(B1:B4)} addiert alle Zellen von \lsg{B1} bis \lsg{B4}.

\vspace{0.3em}
Der \textbf{Doppelpunkt} \texttt{:} bedeutet: \lsg{„von ... bis"}
\end{lueckenkasten}

% ═══════════════════════════════════════════════════════════════
% KASTEN 6: Referenzprinzip
% ═══════════════════════════════════════════════════════════════

\textbf{\textcolor{fachfarbe}{Kasten 6: Das Referenzprinzip}}

\begin{lueckenkasten}
\textbf{Referenzprinzip:}

Ändert sich der Wert einer Zelle, werden alle Formeln, die diese Zelle verwenden,

\lsg{automatisch} neu berechnet.
\end{lueckenkasten}

\vspace{0.3em}
\hrule
\vspace{0.5em}

% ═══════════════════════════════════════════════════════════════
% AUFGABE 1: Zelladressen bestimmen
% ═══════════════════════════════════════════════════════════════

\textbf{\textcolor{fachfarbe}{Aufgabe 1: Zelladressen bestimmen}}

Bestimme die Adresse der markierten Zellen.

\begin{center}
\begin{tikzpicture}[scale=0.7]
    % Header
    \fill[gray!30] (0,0) rectangle (1.2,0.7);
    \fill[gray!30] (1.2,0) rectangle (3.2,0.7);
    \fill[gray!30] (3.2,0) rectangle (5.2,0.7);
    \fill[gray!30] (5.2,0) rectangle (7.2,0.7);
    \draw (0,0) rectangle (1.2,0.7);
    \draw (1.2,0) rectangle (3.2,0.7);
    \draw (3.2,0) rectangle (5.2,0.7);
    \draw (5.2,0) rectangle (7.2,0.7);
    \node at (0.6,0.35) {\small};
    \node at (2.2,0.35) {\small\textbf{A}};
    \node at (4.2,0.35) {\small\textbf{B}};
    \node at (6.2,0.35) {\small\textbf{C}};

    % Row 1
    \fill[gray!30] (0,-0.7) rectangle (1.2,0);
    \draw (0,-0.7) rectangle (1.2,0);
    \draw (1.2,-0.7) rectangle (3.2,0);
    \draw[fachfarbe, thick] (3.2,-0.7) rectangle (5.2,0);
    \draw (5.2,-0.7) rectangle (7.2,0);
    \node at (0.6,-0.35) {\small\textbf{1}};
    \node at (2.2,-0.35) {\small Brezel};
    \node at (4.2,-0.35) {\small 0,80};
    \node at (6.2,-0.35) {\small 12};

    % Row 2
    \fill[gray!30] (0,-1.4) rectangle (1.2,-0.7);
    \draw (0,-1.4) rectangle (1.2,-0.7);
    \draw (1.2,-1.4) rectangle (3.2,-0.7);
    \draw (3.2,-1.4) rectangle (5.2,-0.7);
    \draw[fachfarbe, thick] (5.2,-1.4) rectangle (7.2,-0.7);
    \node at (0.6,-1.05) {\small\textbf{2}};
    \node at (2.2,-1.05) {\small Apfel};
    \node at (4.2,-1.05) {\small 0,50};
    \node at (6.2,-1.05) {\small 8};

    % Row 3
    \fill[gray!30] (0,-2.1) rectangle (1.2,-1.4);
    \draw (0,-2.1) rectangle (1.2,-1.4);
    \draw[fachfarbe, thick] (1.2,-2.1) rectangle (3.2,-1.4);
    \draw (3.2,-2.1) rectangle (5.2,-1.4);
    \draw (5.2,-2.1) rectangle (7.2,-1.4);
    \node at (0.6,-1.75) {\small\textbf{3}};
    \node at (2.2,-1.75) {\small Saft};
    \node at (4.2,-1.75) {\small 1,20};
    \node at (6.2,-1.75) {\small 5};
\end{tikzpicture}
\end{center}

\begin{tabular}{@{}p{4cm}p{4cm}p{4cm}@{}}
a) Zelle mit 0,80: \lsg{B1} & b) Zelle mit 8: \lsg{C2} & c) Zelle mit Saft: \lsg{A3}
\end{tabular}

\vspace{0.3em}

% ═══════════════════════════════════════════════════════════════
% AUFGABE 2: Formeln verstehen
% ═══════════════════════════════════════════════════════════════

\textbf{\textcolor{fachfarbe}{Aufgabe 2: Formeln verstehen}}

Was berechnen diese Formeln? Kreuze an.

\begin{tabular}{@{}p{3.5cm}p{11cm}@{}}
a) \texttt{=A1+B1} & \lsg{$\boxtimes$} addiert A1 und B1 \quad $\square$ multipliziert A1 und B1 \\[0.5em]
b) \texttt{=C2*D2} & $\square$ addiert C2 und D2 \quad \lsg{$\boxtimes$} multipliziert C2 und D2 \\[0.5em]
c) \texttt{=SUMME(B1:B5)} & $\square$ addiert nur B1 und B5 \quad \lsg{$\boxtimes$} addiert B1 bis B5 (alle)
\end{tabular}

\vspace{0.3em}

% ═══════════════════════════════════════════════════════════════
% AUFGABE 3: Fehler finden
% ═══════════════════════════════════════════════════════════════

\textbf{\textcolor{fachfarbe}{Aufgabe 3: Fehler finden}}

Leon hat Formeln eingegeben, aber es funktioniert nicht. Was ist falsch?

\begin{tabular}{@{}p{4cm}p{10cm}@{}}
a) \texttt{B2*C2} & Fehler: \lsg{= am Anfang fehlt} \\[0.5em]
b) \texttt{=SUMME B1:B4} & Fehler: \lsg{Klammern ( ) fehlen}
\end{tabular}

\vspace{0.3em}

% ═══════════════════════════════════════════════════════════════
% KASTEN 7: Aufziehen (Fill-Handle)
% ═══════════════════════════════════════════════════════════════

\textbf{\textcolor{fachfarbe}{Kasten 7: Aufziehen (Fill-Handle)}}

\begin{lueckenkasten}
Mit dem \textbf{Aufzieh-Trick} kann man Formeln automatisch \lsg{kopieren}.

\vspace{0.3em}
Der \textbf{grüne Punkt} befindet sich \lsg{rechts unten} in der Zelle.

\vspace{0.3em}
Beim Aufziehen wird die Formel automatisch angepasst:

\texttt{=B2*C2} $\rightarrow$ \texttt{=B\lsg{3}*C\lsg{3}} $\rightarrow$ \texttt{=B\lsg{4}*C\lsg{4}}
\end{lueckenkasten}

\vspace{0.3em}

% ═══════════════════════════════════════════════════════════════
% AUFGABE 4: Schrittfolge Aufziehen
% ═══════════════════════════════════════════════════════════════

\textbf{\textcolor{fachfarbe}{Aufgabe 4: Schrittfolge Aufziehen}}

Bringe die Schritte in die richtige Reihenfolge (1--4):

\begin{tabular}{@{}c p{12cm}@{}}
\lsg{\fbox{4}} & Den grünen Punkt nach unten ziehen (z.B. bis D6) \\[0.4em]
\lsg{\fbox{1}} & Die erste Formel eingeben (z.B. \texttt{=B2*C2} in D2) \\[0.4em]
\lsg{\fbox{3}} & Die Zelle mit der Formel anklicken (z.B. D2) \\[0.4em]
\lsg{\fbox{2}} & Enter drücken
\end{tabular}

\vspace{0.5em}

% ═══════════════════════════════════════════════════════════════
% AUFGABE 5: Aufziehen verstehen
% ═══════════════════════════════════════════════════════════════

\textbf{\textcolor{fachfarbe}{Aufgabe 5: Aufziehen verstehen}}

In D2 steht die Formel \texttt{=B2*C2}. Du ziehst die Formel nach unten.

\vspace{0.3em}
\begin{tabular}{@{}p{7cm}p{7cm}@{}}
a) Was steht dann in D3? & \texttt{=}\lsg{B3*C3} \\[0.5em]
b) Was steht dann in D5? & \texttt{=}\lsg{B5*C5}
\end{tabular}

\vspace{0.5em}

% ═══════════════════════════════════════════════════════════════
% AUFGABE 6: Autovervollständigung
% ═══════════════════════════════════════════════════════════════

\textbf{\textcolor{fachfarbe}{Aufgabe 6: Autovervollständigung}}

Du möchtest die SUMME-Funktion nutzen. Welche Buchstaben tippst du ein,

damit die \textbf{Autovervollständigung} dir \texttt{SUMME} vorschlägt?

\vspace{0.3em}
Antwort: \texttt{=}\lsg{SU}

\vspace{0.5em}

% ═══════════════════════════════════════════════════════════════
% ZUSAMMENFASSUNG (grau)
% ═══════════════════════════════════════════════════════════════

\begin{merkkasten}
\textbf{Zusammenfassung:} Eine \textbf{Zelle} hat eine Adresse (z.B. B2). \textbf{Formeln} beginnen mit \texttt{=}. Die \textbf{SUMME}-Funktion addiert einen Bereich. Das \textbf{Referenzprinzip} sorgt für automatische Neuberechnung. Mit dem \textbf{Aufziehen} kopierst du Formeln schnell nach unten.
\end{merkkasten}

\end{document}
